\subsection{Numerical Linear Algebra}

% **Role:** quantitative matrix behavior and computation.
%
% Direct solves, QR algorithms, SVD numerics, least-squares stability, sparse/banded structure, Krylov methods, iterative methods, preconditioning preview.

Numerical linear algebra turns matrix structure into algorithms for direct solves, least squares, eigenvalues, SVD, structured transforms, Krylov methods, iteration, and preconditioning.

\begin{perspective} \label{per:solve-with-factors-not-inverses} \hlt{Solve with Factors Before Forming Inverses}\\
For $A\mathbf{x}=\mathbf{b}$, use available factors such as LU, QR, or Cholesky to solve triangular or orthogonal subproblems.\\
This avoids forming $A^{-1}$ and preserves sparse or triangular structure.
\end{perspective}

\begin{perspective} \label{per:iterative-refinement-residual-correction} \hlt{Iterative Refinement by Residual Correction}\\
Given an approximate solution $\mathbf{x}_k$ of $A\mathbf{x}=\mathbf{b}$, set $\mathbf{r}_k=\mathbf{b}-A\mathbf{x}_k$.\\
Solve $A\mathbf{d}_k=\mathbf{r}_k$ with the same factors or preconditioner, then update $\mathbf{x}_{k+1}=\mathbf{x}_k+\mathbf{d}_k$.\\
Exact arithmetic corrects in one step; floating point helps when $\mathbf{r}_k$ is accurate and conditioning is moderate.
\end{perspective}

\begin{method} \label{meth:augmented-inverse-computation-preview} \hlt{Augmented Inverse Computation Preview}\\
For $A\in\operatorname{Mat}(n,n,\mathbb{F})$:
\begin{enumerate}[label=(\roman*)]
\setlength{\itemsep}{0pt}
\item Form $[\,A\ I_n\,]$.
\item Apply Gauss--Jordan reduction to the left block.
\item If $[\,A\ I_n\,]\sim[\,I_n\ B\,]$, then $B=A^{-1}$.
\item Otherwise $A$ is not invertible, by Corollary \ref{cor:square-matrix-invertible-iff-full-rank}.
\end{enumerate}
\end{method}

\begin{proposition} \label{prop:qr-least-squares-and-projection} \hlt{QR Route for Least Squares}\\
Let $A=QR\in\operatorname{Mat}(m,n,\mathbb{F})$ be a full-column-rank narrow QR factorisation. For $\mathbf{y}\in\mathbb{F}^m$, the least-squares solution solves $R\mathbf{x}=Q^*\mathbf{y}$. Moreover $P_{\operatorname{col}(A)}=QQ^*$.
\begin{proof}
The normal equations reduce by $A=QR$ and $Q^*Q=I_n$ to $R^*R\mathbf{x}=R^*Q^*\mathbf{y}$.
Since $R$ is invertible, this is equivalent to $R\mathbf{x}=Q^*\mathbf{y}$. Also $\operatorname{col}(A)=\operatorname{col}(Q)$, and the columns of $Q$ are orthonormal, so $P_{\operatorname{col}(A)}=QQ^*$.\\
\hlt{Structural Commentary}: QR replaces $\operatorname{col}(A)$ by an orthonormal basis $Q$.\\
\hlt{Key Insight}: Project onto $Q$, then solve triangularly.
\end{proof}
\end{proposition}

\begin{perspective} \label{per:normal-equations-squaring-condition-number} \hlt{Why Normal Equations Can Be Numerically Fragile}\\
For full-column-rank least squares, solving $A^*A\mathbf{x}=A^*\mathbf{y}$ squares the spectral condition number because $\kappa_2(A^*A)=\kappa_2(A)^2$. QR-based methods avoid this squaring and therefore preserve conditioning better.
\end{perspective}

\begin{figure}[H]
  \centering
  \begin{tikzpicture}[
  font=\small,
  cell/.style={minimum width=0.48cm, minimum height=0.40cm, inner sep=1pt, outer sep=0pt, anchor=center},
  aband/.style={fill=gray!16},
  lband/.style={fill=blue!10, text=blue!70!black},
  uband/.style={fill=gray!16},
  title/.style={font=\small}
]
  \matrix (A) [
    matrix of math nodes,
    nodes={cell},
    left delimiter={[},
    right delimiter={]},
    row sep=0pt,
    column sep=0pt
  ] {
    |[aband]| * & |[aband]| * & \phantom{*} & \phantom{*} & \phantom{*}\\
    |[aband]| * & |[aband]| * & |[aband]| * & \phantom{*} & \phantom{*}\\
    \phantom{*} & |[aband]| * & |[aband]| * & |[aband]| * & \phantom{*}\\
    \phantom{*} & \phantom{*} & |[aband]| * & |[aband]| * & |[aband]| *\\
    \phantom{*} & \phantom{*} & \phantom{*} & |[aband]| * & |[aband]| *\\
  };
  \node[title, above=5pt] at (A.north) {$A$};

  \node at (2.35,0) {$=$};

  \begin{scope}[shift={(4.35,0)}]
  \matrix (L) [
    matrix of math nodes,
    nodes={cell},
    left delimiter={[},
    right delimiter={]},
    row sep=0pt,
    column sep=0pt
  ] {
    |[lband]| 1 & \phantom{*} & \phantom{*} & \phantom{*} & \phantom{*}\\
    |[lband]| * & |[lband]| 1 & \phantom{*} & \phantom{*} & \phantom{*}\\
    \phantom{*} & |[lband]| * & |[lband]| 1 & \phantom{*} & \phantom{*}\\
    \phantom{*} & \phantom{*} & |[lband]| * & |[lband]| 1 & \phantom{*}\\
    \phantom{*} & \phantom{*} & \phantom{*} & |[lband]| * & |[lband]| 1\\
  };
  \node[title, above=5pt] at (L.north) {$L$};
  \end{scope}

  \node at (6.55,0) {$\cdot$};

  \begin{scope}[shift={(8.55,0)}]
  \matrix (U) [
    matrix of math nodes,
    nodes={cell},
    left delimiter={[},
    right delimiter={]},
    row sep=0pt,
    column sep=0pt
  ] {
    |[uband]| * & |[uband]| * & \phantom{*} & \phantom{*} & \phantom{*}\\
    \phantom{*} & |[uband]| * & |[uband]| * & \phantom{*} & \phantom{*}\\
    \phantom{*} & \phantom{*} & |[uband]| * & |[uband]| * & \phantom{*}\\
    \phantom{*} & \phantom{*} & \phantom{*} & |[uband]| * & |[uband]| *\\
    \phantom{*} & \phantom{*} & \phantom{*} & \phantom{*} & |[uband]| *\\
  };
  \node[title, above=5pt] at (U.north) {$U$};
  \end{scope}
\end{tikzpicture}

  \caption{A tridiagonal matrix can factor as bidiagonal $L$ times bidiagonal $U$, retaining local band structure.}
  \label{fig:band-matrix-lu-structure}
\end{figure}

\begin{definition} \label{def:band-and-tridiagonal-matrix} \hlt{Band and Tridiagonal Matrices}\\
For $A=[a_{ij}]\in\operatorname{Mat}(n,n,\mathbb{F})$ and $1\leq w\leq n$, say that $A$ has \hlt{half-bandwidth} $w$ if $a_{ij}=0$ whenever $\abs{i-j}\geq w$. Thus $w=1$ gives a diagonal matrix, while $w=2$ gives a \hlt{tridiagonal} matrix.
\end{definition}

\begin{method} \label{meth:banded-elimination-preview} \hlt{Banded Elimination Preview}\\
Let $A\in\operatorname{Mat}(n,n,\mathbb{F})$ have half-bandwidth $w$.
\begin{enumerate}[label=(\roman*)]
\setlength{\itemsep}{0pt}
\item Eliminate without row exchanges that leave the band.
\item At pivot $k$, update only entries with $\abs{i-k}<w$ and $\abs{j-k}<w$.
\item Obtain $A=LU$ with $L$ and $U$ of comparable bandwidth.
\end{enumerate}
Cost: $O(nw^2)$ operations; for tridiagonal matrices $(w=2)$, $O(n)$.
\end{method}

\begin{perspective} \label{per:five-point-stencil-bandwidth} \hlt{Five-Point Stencils Are Banded but Less Local}\\
The $1$D second-difference pattern $-1,2,-1$ is tridiagonal; the row-ordered $N\times N$ five-point stencil is still sparse, but vertical neighbors create diagonals offset by about $N$. Thus bandwidth grows with grid width, and incomplete LU keeps a prescribed sparse pattern for a cheap approximate solve or preconditioner.
\end{perspective}

\begin{figure}[H]
  \centering
  \begin{tikzpicture}[
  >=Stealth,
  every node/.style={font=\small},
  box/.style={draw=gray!50, fill=gray!3, rounded corners=2pt, line width=0.55pt},
  main/.style={line width=1.0pt, black},
  near/.style={line width=0.8pt, blue!70!black},
  far/.style={line width=0.8pt, teal!70!black},
  bracket/.style={gray!75!black, line width=0.55pt},
  legend/.style={font=\footnotesize, anchor=west}
]
  \node at (-0.62,0) {$A=$};

  \draw[bracket] (-0.18,1.32) -- (-0.34,1.32) -- (-0.34,-1.32) -- (-0.18,-1.32);
  \draw[bracket] (3.18,1.32) -- (3.34,1.32) -- (3.34,-1.32) -- (3.18,-1.32);

  \draw[box] (0,-1.2) rectangle (3,1.2);
  \begin{scope}
    \clip (0,-1.2) rectangle (3,1.2);
    \draw[far] (-0.1,0.45) -- (1.75,-1.35);
    \draw[near] (-0.1,0.95) -- (2.45,-1.35);
    \draw[main] (0.25,1.25) -- (2.9,-1.35);
    \draw[near] (0.75,1.35) -- (3.15,-0.95);
    \draw[far] (1.55,1.35) -- (3.15,-0.20);
  \end{scope}

  \draw[main] (3.9,0.72) -- (4.35,0.72);
  \node[legend] at (4.45,0.72) {offset $0$};
  \draw[near] (3.9,0.12) -- (4.35,0.12);
  \node[legend, blue!70!black] at (4.45,0.12) {offset $\pm1$};
  \draw[far] (3.9,-0.48) -- (4.35,-0.48);
  \node[legend, teal!70!black] at (4.45,-0.48) {offset $\pm N$};
\end{tikzpicture}

  \caption{Row ordering gives matrix bands at offsets $0,\pm1,\pm N$.}
  \label{fig:five-point-stencil-bandwidth}
\end{figure}

\begin{definition} \label{def:upper-hessenberg-matrix} \hlt{Upper Hessenberg Matrix}\\
Let $A=[a_{ij}]\in\operatorname{Mat}(n,n,\mathbb{F})$. Then $A$ is \hlt{upper Hessenberg} if $a_{ij}=0$ whenever $i>j+1$.\\
Equivalently, all entries below the first subdiagonal are zero.
\end{definition}

Numerical linear algebra also replaces a matrix by a unitarily similar one with many forced zero entries. Upper Hessenberg form is the standard preliminary reduction before QR-type eigenvalue algorithms.

\begin{figure}[H]
  \centering
  \begin{tikzpicture}[font=\small]
  \node at (0,0) {$
  \begin{bmatrix}
  * & * & * & * & *\\
  * & * & * & * & *\\
  0 & * & * & * & *\\
  0 & 0 & * & * & *\\
  0 & 0 & 0 & * & *
  \end{bmatrix}
  $};
  \draw[red!70, thick] (-1.35,-0.45) -- (0.65,-1.15);
  \node[red!70!black, align=center] at (2.45,-1.0) {first\\subdiagonal};
  \draw[->, red!70, thick] (1.65,-0.95) -- (0.72,-0.92);

  \node[blue!70!black, align=center] at (0,1.75) {entries below this subdiagonal are zero};
  \draw[->, blue!70, thick] (-0.2,1.45) -- (-1.0,-0.85);
\end{tikzpicture}

  \caption{Upper Hessenberg pattern: red is the first subdiagonal; blue entries below it are forced zero.}
  \label{fig:upper-hessenberg-pattern}
\end{figure}

\begin{theorem} \label{thm:unitary-upper-hessenberg-reduction} \hlt{Unitary Upper Hessenberg Reduction}\\
Every $A\in\operatorname{Mat}(n,n,\mathbb{F})$ is unitarily similar to upper Hessenberg matrix with real nonnegative subdiagonal entries.
\begin{proof}
Set $A_0=A$. For $k=1,\ldots,n-2$, suppose columns $1,\ldots,k-1$ of $A_{k-1}$ already have zeros below the first subdiagonal. Let $\mathbf{w}$ be the tail of column $k$ in rows $k+1,\ldots,n$.\\
By Corollary \ref{cor:householder-unitary-annihilator}, choose a unitary $V_k$ with $V_k^*\mathbf{w}=\alpha\mathbf{e}_1$, $\alpha\geq0$.\\
Put $U_k=I_k\oplus V_k$ and $A_k=U_k^*A_{k-1}U_k$. Right multiplication by $U_k$ changes only columns $k+1,\ldots,n$, so columns $1,\ldots,k-1$ keep their existing Hessenberg zeros. Left multiplication by $U_k^*$ acts only on rows $k+1,\ldots,n$; those rows are already zero in columns $1,\ldots,k-1$, so the zeros remain. In column $k$, the tail becomes $V_k^*\mathbf{w}=\alpha\mathbf{e}_1$, so entries below row $k+1$ are zero and the new subdiagonal entry is real nonnegative.\\
After $n-2$ steps, $A_{n-2}=U^*AU$ is upper Hessenberg, where $U=U_1\cdots U_{n-2}$ is unitary.\\
\hlt{Structural Commentary}: Each trailing unitary creates one Hessenberg column and preserves earlier ones.\\
\hlt{Key Insight}: Kill one column tail at a time.
\end{proof}
\end{theorem}

\begin{perspective} \label{per:hessenberg-before-qr} \hlt{Why Hessenberg Form Matters}\\
Hessenberg form preserves spectral data by unitary similarity, making QR-type eigenvalue algorithms cheaper.
\end{perspective}

\begin{method} \label{meth:qr-eigenvalue-algorithm-preview} \hlt{QR Eigenvalue Algorithm Preview}\\
For a square matrix $A_0$, the basic QR eigenvalue iteration is:
\begin{enumerate}[label=(\roman*)]
\setlength{\itemsep}{0pt}
\item In practice, first reduce $A_0$ to Hessenberg form by Theorem \ref{thm:unitary-upper-hessenberg-reduction}.
\item Factor $A_k=Q_kR_k$, with $Q_k$ unitary and $R_k$ upper triangular.
\item Set $A_{k+1}=R_kQ_k=Q_k^*A_kQ_k$, using shifts to speed convergence.
\end{enumerate}
Each step is a unitary similarity, so the spectrum is preserved.
\end{method}

\begin{method} \label{meth:hessenberg-qr-step-plane-rotations} \hlt{Hessenberg QR Step by Plane Rotations}\\
Let $H\in\operatorname{Mat}(n,n,\mathbb{C})$ be upper Hessenberg and let $\mu\in\mathbb{C}$.
\begin{enumerate}[label=(\roman*)]
\setlength{\itemsep}{0pt}
\item Set $B_0=H-\mu I_n$.
\item For $j=1,\ldots,n-1$, choose a unitary plane rotation $G_j$ acting on coordinates $j,j+1$,\\
so that $(G_j^*B_{j-1})_{j+1,j}=0$ and $B_j=G_j^*B_{j-1}$.
\item Put $Q=G_1\cdots G_{n-1}$ and $R=B_{n-1}$; then $H-\mu I_n=QR$.
\item Define the shifted QR update $H_+=RQ+\mu I_n=Q^*HQ$
\end{enumerate}
Adjacent rotations implement the Hessenberg step by local bulge chasing.
\end{method}

\begin{theorem} \label{thm:power-method-dominant-eigenvector-preview} \hlt{Power Method Dominant-Eigenvector Preview}\\
Let $A\mathbf{v}_j=\lambda_j\mathbf{v}_j$, where $\mathbf{v}_1,\ldots,\mathbf{v}_n$ is a basis of $\mathbb{C}^n$, and assume $\abs{\lambda_1}>\max_{j\geq2}\abs{\lambda_j}$.\\
If $\mathbf{u}_0=\sum\limits_{j=1}^{n}c_j\mathbf{v}_j$ with $c_1\neq0$, then $\lambda_1^{-k}A^k\mathbf{u}_0\longrightarrow c_1\mathbf{v}_1$. Thus the directions of $A^k\mathbf{u}_0$ converge to $\operatorname{span}\{\mathbf{v}_1\}$.
\begin{proof}
Set $\rho=\max\limits_{j\geq2}\abs{\lambda_j/\lambda_1}<1$. Since $A^k\mathbf{v}_j=\lambda_j^k\mathbf{v}_j$, $\lambda_1^{-k}A^k\mathbf{u}_0 =c_1\mathbf{v}_1+\sum\limits_{j=2}^{n}c_j\left(\frac{\lambda_j}{\lambda_1}\right)^k\mathbf{v}_j$.\\
The remainder has norm at most $\rho^k\sum\limits_{j=2}^{n}\abs{c_j}\norm{\mathbf{v}_j}$, hence tends to $0$.\\
\hlt{Structural Commentary}: The power method is spectral filtering by repeated multiplication.\\
\hlt{Key Insight}: A spectral gap makes every non-dominant component decay relative to the dominant one.
\end{proof}
\end{theorem}

\begin{method} \label{meth:inverse-and-shifted-power-iteration-preview} \hlt{Inverse and Shifted Power Iteration Preview}\\
Let $\mathcal{N}(\mathbf{z})\defeq\mathbf{z}/\norm{\mathbf{z}}_2$ for $\mathbf{z}\neq\mathbf{0}$.
\begin{enumerate}[label=(\roman*)]
\setlength{\itemsep}{0pt}
\item \hlt{Inverse}: if $0\notin\operatorname{spec}(A)$, solve $A\mathbf{y}_{k+1}=\mathbf{x}_k$ and set $\mathbf{x}_{k+1}=\mathcal{N}(\mathbf{y}_{k+1})$.\\
This is power iteration on $A^{-1}$ and targets minimal $\abs{\lambda}$.
\item \hlt{Shift-invert}: for $\sigma\notin\operatorname{spec}(A)$, $(A-\sigma I)\mathbf{y}_{k+1}=\mathbf{x}_k$, $\mathbf{x}_{k+1}=\mathcal{N}(\mathbf{y}_{k+1})$.\\
Eigenvalues become $(\lambda-\sigma)^{-1}$, targeting $\lambda$ nearest $\sigma$.
\item \hlt{Shifted power}: $\mathbf{x}_{k+1}=\mathcal{N}\bigl((A-\sigma I)\mathbf{x}_k\bigr)$; eigenvalues become $\lambda-\sigma$, emphasizing large $\abs{\lambda-\sigma}$.
\end{enumerate}
Use factors of $A$ or $A-\sigma I$ for solves; do not form explicit inverses.
\end{method}

\begin{method} \label{meth:svd-numerics-preview} \hlt{SVD Numerics Preview}\\
Stable numerical SVD computation proceeds as:
\begin{enumerate}[label=(\roman*)]
\setlength{\itemsep}{0pt}
\item Avoid forming $A^*A$ when possible, since it can square the condition number.
\item Use orthogonal or unitary transformations to reduce $A$ to bidiagonal form.
\item Compute singular values from the bidiagonal problem, then use them for low-rank approximation, pseudoinverse truncation, and least-squares diagnostics through Theorem \ref{thm:eckart-young-frobenius-spectral} and Definition \ref{def:spectral-condition-number}.
\end{enumerate}
\end{method}

\begin{definition} \label{def:discrete-fourier-matrix} \hlt{Discrete Fourier Matrix}\\
For $n\geq1$, set $\omega_n=e^{2\pi i/n}$ and $F_n\defeq[\omega_n^{jk}]_{j,k=0}^{n-1}$.\\
Then $F_n\mathbf{c}$ is the unnormalised DFT of $\mathbf{c}\in\C^n$; $Q_n\defeq n^{-1/2}F_n$ is the normalized Fourier matrix.
\end{definition}

\begin{proposition} \label{prop:fourier-matrix-inverse} \hlt{Inverse of the Fourier Matrix}\\
The Fourier matrix satisfies $F_n^*F_n=nI_n$, and therefore $F_n^{-1}=\frac1nF_n^*$.
\begin{proof}
Theorem \ref{thm:fourier-basis-diagonalizes-cyclic-circulants} gives $Q_n^*Q_n=I_n$ for $Q_n=n^{-1/2}F_n$. Multiplying by $n$ gives $F_n^*F_n=nI_n$, so $n^{-1}F_n^*$ is a left inverse of $F_n$. Since $F_n$ is square, Theorem \ref{thm:square-one-sided-inverse-is-two-sided} makes it the inverse.\\
\hlt{Structural Commentary}: Fourier inversion is conjugate transpose plus one scalar normalization.\\
\hlt{Key Insight}: The inverse is simple because the Fourier columns are orthogonal.
\end{proof}
\end{proposition}

\begin{method} \label{meth:radix-two-fft-split} \hlt{Radix-Two FFT Split}\\
Let $n=2m$ and $\mathbf{y}=F_n\mathbf{c}$.
\begin{enumerate}[label=(\roman*)]
\setlength{\itemsep}{0pt}
\item Split $\mathbf{c}'\defeq(c_{2r})_{r=0}^{m-1}$ and $\mathbf{c}''\defeq(c_{2r+1})_{r=0}^{m-1}$.
\item Compute $\mathbf{y}'\defeq F_m\mathbf{c}'$ and $\mathbf{y}''\defeq F_m\mathbf{c}''$.
\item With $d_j\defeq\omega_n^j$, and using $\omega_n^2=\omega_m$ and $\omega_n^m=-1$, set $y_j=y'_j+d_jy''_j$; and $y_{j+m}=y'_j-d_jy''_j$ for $0\leq j<m$.
\item For $n=2^\ell$, repeat (i)--(iii) for $\ell$ stages.\\
$n^2$ dense products become about $\frac12 n\log_2 n$ twiddle multiplications plus additions.
\end{enumerate}
\end{method}

\begin{figure}[H]
  \centering
  \begin{tikzpicture}[
  >=Stealth,
  every node/.style={font=\small},
  vector/.style={draw=gray!60, fill=white, rounded corners=2pt, minimum width=1.25cm, minimum height=0.58cm, align=center},
  block/.style={draw=gray!65, fill=gray!8, rounded corners=3pt, minimum width=1.25cm, minimum height=0.62cm, align=center},
  map/.style={->, line width=0.9pt, blue!70!black}
]
  \node[vector] (c) at (0,0) {$\mathbf{c}$};
  \node[block] (split) at (1.85,0) {split};

  \node[vector] (ce) at (3.75,0.8) {$\mathbf{c}'=(c_0,c_2)^t$};
  \node[vector] (co) at (3.75,-0.8) {$\mathbf{c}''=(c_1,c_3)^t$};

  \node[block] (f2e) at (5.8,0.8) {$F_2$};
  \node[block] (f2o) at (5.8,-0.8) {$F_2$};

  \node[vector] (ye) at (7.55,0.8) {$\mathbf{y}'$};
  \node[vector] (yo) at (7.55,-0.8) {$\mathbf{y}''$};

  \node[block, minimum width=1.45cm, minimum height=0.82cm] (butterfly) at (9.75,0) {$B_4$};
  \node[vector] (y) at (11.85,0) {$\mathbf{y}=F_4\mathbf{c}$};

  \draw[map] (c) -- (split);
  \draw[map] (split.east) -- (ce.west);
  \draw[map] (split.east) -- (co.west);
  \draw[map] (ce) -- (f2e);
  \draw[map] (co) -- (f2o);
  \draw[map] (f2e) -- (ye);
  \draw[map] (f2o) -- (yo);
  \draw[map] (ye.east) -- (butterfly.west);
  \draw[map] (yo.east) -- (butterfly.west);
  \draw[map] (butterfly) -- (y);

\end{tikzpicture}

  \caption{For $n=4$: split, apply two $F_2$ transforms, then $B_4\binom{\mathbf{y}'}{\mathbf{y}''}=\binom{\mathbf{y}'+D\mathbf{y}''}{\mathbf{y}'-D\mathbf{y}''}$ with $D=\operatorname{diag}(1,i)$.}
  \label{fig:fft-butterfly-n4}
\end{figure}

\begin{definition} \label{def:krylov-subspace} \hlt{Krylov Subspace}\\
For $A\in\operatorname{Mat}(n,n,\mathbb{F})$ and a vector $\mathbf{b}\in\mathbb{F}^n$, the $m$th \hlt{Krylov subspace} generated by $(A,\mathbf{b})$ is\\
$\mathcal{K}_m(A,\mathbf{b})\defeq\operatorname{span}\{\mathbf{b},A\mathbf{b},\ldots,A^{m-1}\mathbf{b}\}$.
\end{definition}

\begin{perspective} \label{per:krylov-methods-preview} \hlt{Krylov Methods Preview}\\
Krylov methods work in $\mathcal{K}_m(A,\mathbf{b})$: Arnoldi and Lanczos build bases, while CG and GMRES impose projection or residual-minimization conditions using only matrix-vector products.
\end{perspective}

\begin{method} \label{meth:arnoldi-lanczos-cg-preview} \hlt{Arnoldi, Lanczos, and CG Preview}\\
Let $\mathcal{K}_m=\mathcal{K}_m(A,\mathbf{b})$ and let $Q_m=[\mathbf{q}_1\ \cdots\ \mathbf{q}_m]$ be an orthonormal basis.
\begin{enumerate}[label=(\roman*)]
\setlength{\itemsep}{0pt}
\item \hlt{Arnoldi}: $AQ_m=Q_mH_m+h_{m+1,m}\mathbf{q}_{m+1}\mathbf{e}_m^t$, with $H_m$ upper Hessenberg.
\item \hlt{Lanczos}: if $A=A^*$, then $H_m$ is tridiagonal and $A\mathbf{q}_j=\beta_{j-1}\mathbf{q}_{j-1}+\alpha_j\mathbf{q}_j+\beta_j\mathbf{q}_{j+1}$.
\item \hlt{CG}: if $A\succ0$, choose $\mathbf{x}_k\in\mathbf{x}_0+\mathcal{K}_k(A,\mathbf{r}_0)$ with $\mathbf{r}_k\perp\mathcal{K}_k(A,\mathbf{r}_0)$; search directions satisfy $\mathbf{p}_i^*A\mathbf{p}_j=0$, $i\neq j$.
\end{enumerate}
Each step uses one matrix-vector product.
\end{method}

\begin{definition} \label{def:matrix-splitting-iteration} \hlt{Matrix Splitting Iteration}\\
Let $A=S-T$ with $S$ invertible, and set $B\defeq S^{-1}T$.\\
The \hlt{splitting iteration} for $A\mathbf{x}=\mathbf{b}$ is $S\mathbf{x}_{k+1}=T\mathbf{x}_k+\mathbf{b}$, and $\mathbf{x}_{k+1}=B\mathbf{x}_k+S^{-1}\mathbf{b}$.\\
For a solution $\mathbf{x}_*$, the error $\mathbf{e}_k\defeq\mathbf{x}_k-\mathbf{x}_*$ satisfies $\mathbf{e}_{k+1}=B\mathbf{e}_k$.
\end{definition}

\begin{theorem} \label{thm:splitting-iteration-convergence-spectral-radius} \hlt{Splitting Convergence Criterion}\\
Let $A=S-T$, with $S$ invertible, and set $B=S^{-1}T$. For $\mathbf{x}_{k+1}=B\mathbf{x}_k+S^{-1}\mathbf{b}$:
\begin{enumerate}[label=(\roman*)]
\setlength{\itemsep}{0pt}
\item if $A$ is invertible, then $\mathbf{x}_k\to A^{-1}\mathbf{b}$ for all $\mathbf{x}_0,\mathbf{b}$ if and only if $\rho(B)<1$;
\item $\rho(B)<1$ implies $A$ is invertible.
\end{enumerate}
\begin{proof}
\begin{enumerate}[label=(\roman*)]
\setlength{\itemsep}{0pt}
\item Let $\mathbf{x}_*=A^{-1}\mathbf{b}$ and $\mathbf{e}_k=\mathbf{x}_k-\mathbf{x}_*$. Then $\mathbf{e}_k=B^k\mathbf{e}_0$, so convergence for all $\mathbf{x}_0,\mathbf{b}$ is equivalent to $B^k\to0$, equivalently $\rho(B)<1$ by Theorem \ref{thm:matrix-powers-converge-to-zero-spectral-radius}.
\item If $\rho(B)<1$, then $1\notin\operatorname{spec}(B)$, so $I_n-B$ is invertible. Since $A=S-T=S(I_n-B)$, $A$ is invertible.
\end{enumerate}
\hlt{Structural Commentary}: A stationary method is a matrix-power problem in disguise.\\
\hlt{Key Insight}: The iteration matrix $B=S^{-1}T$ controls the errors.
\end{proof}
\end{theorem}

\begin{corollary} \label{cor:singular-matrix-splitting-obstruction} \hlt{Singular Matrices Obstruct Splitting Convergence}\\
Let $A=S-T$, with $S$ invertible. If $A$ is singular, then $1\in\operatorname{spec}(S^{-1}T)$.\\
Hence $\rho(S^{-1}T)\geq1$, so $\rho(S^{-1}T)<1$ cannot hold.
\begin{proof}
Choose $\mathbf{0}\neq\mathbf{x}\in\ker A$. Then $(S-T)\mathbf{x}=\mathbf{0}$, so $S\mathbf{x}=T\mathbf{x}$ and $S^{-1}T\mathbf{x}=\mathbf{x}$. Thus $1$ is eigenvalue of $S^{-1}T$.\\
\hlt{Structural Commentary}: Singular $A$ forces an undamped error mode.\\
\hlt{Key Insight}: A null vector of $A$ becomes a fixed error mode.
\end{proof}
\end{corollary}

\begin{method} \label{meth:stationary-iterative-method} \hlt{Stationary Iterative Method}\\
For $A\mathbf{x}=\mathbf{b}$, a stationary iteration is:
\begin{enumerate}[label=(\roman*)]
\setlength{\itemsep}{0pt}
\item Choose fixed $M$ and $\mathbf{c}$.
\item Update $\mathbf{x}_{k+1}=M\mathbf{x}_k+\mathbf{c}$.
\item If $\mathbf{x}_*$ exists, use $\mathbf{x}_{k+1}-\mathbf{x}_*=M(\mathbf{x}_k-\mathbf{x}_*)$ and Theorem \ref{thm:matrix-powers-converge-to-zero-spectral-radius} for convergence.
\end{enumerate}
\end{method}

\begin{method} \label{meth:jacobi-gauss-seidel-sor-splittings} \hlt{Jacobi, Gauss--Seidel, and SOR Splittings}\\
Let $A=D+L+U$, with $D$ diagonal, $L$ strictly lower triangular, and $U$ strictly upper triangular. For a splitting $A=S-T$, use $S\mathbf{x}_{k+1}=T\mathbf{x}_k+\mathbf{b}$ and $B=S^{-1}T$.
\begin{enumerate}[label=(\roman*)]
\setlength{\itemsep}{0pt}
\item \hlt{Jacobi}: $S=D$, $T=-(L+U)$, $B_J=-D^{-1}(L+U)$.
\item \hlt{Gauss--Seidel}: $S=D+L$, $T=-U$, $B_{\operatorname{GS}}=-(D+L)^{-1}U$.
\item \hlt{SOR}: for $\omega\neq0$, $S_\omega=\frac1\omega D+L$; $T_\omega=\left(\frac1\omega-1\right)D-U$; and $B_\omega=S_\omega^{-1}T_\omega$.
\end{enumerate}
Convergence is controlled by $\rho(B)$; $\omega=1$ gives Gauss--Seidel.
\end{method}

\begin{corollary} \label{cor:strict-diagonal-dominance-jacobi-convergence} \hlt{Strict Diagonal Dominance Gives Jacobi Convergence}\\
If $A=[a_{ij}]\in\operatorname{Mat}(n,n,\mathbb{C})$ is strictly row diagonally dominant, then Jacobi converges for every right-hand side.
\begin{proof}
Here $a_{ii}\neq0$ and $B_J=-D^{-1}(L+U)=I_n-D^{-1}A$. Thus $(B_J)_{ii}=0$ and $\sum\limits_{j\neq i}\abs{(B_J)_{ij}}=\sum\limits_{j\neq i}\abs{a_{ij}}/\abs{a_{ii}}<1$.\\
Gersgorin gives $\operatorname{spec}(B_J)\subset\{z:\abs{z}<1\}$, so $\rho(B_J)<1$; Theorem \ref{thm:splitting-iteration-convergence-spectral-radius} applies.\\
\hlt{Structural Commentary}: Diagonal dominance puts the Jacobi spectrum in the unit disk.\\
\hlt{Key Insight}: Small off-diagonal row sums force the error matrix to have spectral radius below $1$.
\end{proof}
\end{corollary}

\begin{example} \label{ex:two-by-two-jacobi-gauss-seidel-convergence} \hlt{Two-by-Two Jacobi and Gauss--Seidel}\\
For $A=\begin{bmatrix}2&1\\1&2\end{bmatrix}$, $\mathbf{b}\in\mathbb{R}^2$, write $A=D+L+U$.\\
The Jacobi iteration matrix is $B_J=-D^{-1}(L+U)= \begin{bmatrix} 0&-\frac12\\ -\frac12&0 \end{bmatrix}$, so $\rho(B_J)=\frac12$.\\
The Gauss--Seidel iteration matrix is $B_{\operatorname{GS}}=-(D+L)^{-1}U= \begin{bmatrix} 0&-\frac12\\ 0&\frac14 \end{bmatrix}$, so $\rho(B_{\operatorname{GS}})=\frac14$.\\
Both converge by Theorem \ref{thm:splitting-iteration-convergence-spectral-radius}, and smaller spectral radius predicts faster asymptotic decay for Gauss--Seidel.
\end{example}

\begin{example} \label{ex:poisson-tridiagonal-stationary-rates} \hlt{Poisson Tridiagonal Stationary Rates}\\
For $A_n=\operatorname{tridiag}(-1,2,-1)$, Jacobi gives $B_J=I_n-\frac12A_n$; $\lambda_k(B_J)=\cos\frac{k\pi}{n+1}$; and $\rho(B_J)=\cos\frac{\pi}{n+1}$.\\
For the same ordering, Gauss--Seidel satisfies $\rho(B_{\operatorname{GS}})=\rho(B_J)^2=\cos^2\frac{\pi}{n+1}$.
As $n\to\infty$, both radii tend to $1$; fine grids need acceleration or preconditioning.
\end{example}

\begin{perspective} \label{per:preconditioning-preview} \hlt{Preconditioning Preview}\\
A preconditioner replaces $A\mathbf{x}=\mathbf{b}$ by an equivalent system such as $M^{-1}A\mathbf{x}=M^{-1}\mathbf{b}$ or $AM^{-1}\mathbf{y}=\mathbf{b}$, with $\mathbf{x}=M^{-1}\mathbf{y}$. It should improve conditioning or convergence while being much cheaper to apply than $A^{-1}$.
\end{perspective}
